\documentclass[a4paper,12pt]{article}
% Umlaute im PDF sind nicht zusammengesetzt
\usepackage[T1]{fontenc}
% Umlaute müssen nicht maskiert werden
\usepackage[utf8]{inputenc}
% Titel
\title{Simulation und Realisierung der Gelenkregelung G2 Stäubli TX60L}
% Autor
\author{Fabian Alexander Wilms, Salvador Tazanou}
% Vektorschriftart
\usepackage{lmodern}
% für Deutsche Begriffe und Silbentrennung
\usepackage[ngerman]{babel}
% Formeln
\usepackage{amsmath}
% Grafiken
\usepackage{graphicx}
% Abkürzungsverzeichnis
\usepackage[printonlyused]{acronym}
% Stichwortverzeichnis
\usepackage{makeidx}
\makeindex
% idxlayout works around a bug in LaTeX, which results in a wrong vertical position of the title
% http://tex.stackexchange.com/questions/23287/setting-the-same-distance-from-the-top-of-the-page-
% for-chapter-and-index-titles
\usepackage{idxlayout}
% Zeilenabstand 1,5 Zeilen
\usepackage{setspace}
\onehalfspacing
% BibLaTeX benutzen
\usepackage[backend=biber,citestyle=alphabetic,bibstyle=alphabetic]{biblatex}
\addbibresource{bib/Bibliographie.bib}
% Formatierung an Vorgaben anpassen
\DeclareLabelalphaTemplate{
	\labelelement{
		\field[final]{shorthand}
		\field{label}
		\field[strwidth=3,strside=left,ifnames=1]{labelname}
		\field[strwidth=1,strside=left]{labelname}
	}
	\labelelement{
	\literal{-}
	}
	\labelelement{
	\field[strwidth=2,strside=right]{year}
	}
}
% Seitenränder anpassen
\usepackage{geometry}
\geometry{left=2cm, right=2cm, top=2.5cm, bottom=2cm}
% eps-Grafiken
\usepackage{epstopdf}
% Siehe http://tex.stackexchange.com/questions/28198/using-ieeetrantools
\usepackage[retainorgcmds]{IEEEtrantools}
\usepackage{microtype}
\usepackage{blindtext}
\newcounter{savepage}
% pdf-Metadaten und Hyperlinks
\usepackage[pdfsubject = {{}},
			pdfkeywords = {{Keyword Keyword2}},
			pdfstartview = Fit,
			pdfpagelayout = SinglePage,
			hidelinks,
			pdfusetitle]{hyperref}
% Schneide automatisch alle mit /includegraphics eingebundene PDFs zu
\usepackage{xstring}
\let\oldincludegraphics\includegraphics
\renewcommand{\includegraphics}[2][width=\textwidth]{%
	\immediate\write18{pdfcrop #2}%
	\StrSubstitute{#2}{.pdf}{-crop.pdf}[\temp]%
	\oldincludegraphics[#1]{\temp}%
	}
% Deutsche Anführungszeichen
\usepackage{csquotes}
\begin{document}
\pagenumbering{Roman}
\newpage\null\thispagestyle{empty}\newpage
\newpage
\pagestyle{empty}
\begin{center}
\includegraphics[width=0.35\textwidth]{fig/Logo/h-da-fbmk-logo-sw.pdf}
\vfill
\Large Hochschule Darmstadt \\
\vspace{12pt}
Fachbereich Maschinenbau und Kunststofftechnik \normalsize \\
\vfill
% \@title und \@author verfügbar machen
\makeatletter
\@title \\
\vfill
Dokumentation zum Seminar der Robotik \\
\vfill
vorgelegt von \\
\vspace{12pt}
\@author
\makeatother
\vfill
\begin{tabular}[h]{p{4cm}r}
	Betreuung: & A. König, W. Weber\\
\end{tabular}
\end{center}
\newpage
\pagestyle{plain}
\section*{Eidesstattliche Erklärung}
\addcontentsline{toc}{section}{Eidesstattliche Erklärung}
Ich versichere hiermit, dass ich die vorliegende Arbeit selbständig verfasst und keine anderen als die im Literaturverzeichnis angegebenen Quellen benutzt habe. \\
\noindent Alle Stellen, die wörtlich oder sinngemäß aus veröffentlichten oder noch nicht veröffentlichten Quellen entnommen sind, sind als solche kenntlich gemacht. \\
\noindent Die Zeichnungen oder Abbildungen in dieser Arbeit sind von mir selbst erstellt worden oder mit einem entsprechenden Quellennachweis versehen. \\
\noindent Diese Arbeit ist in gleicher oder ähnlicher Form noch bei keiner anderen Prüfungsbehörde eingereicht worden.
\newline\newline
Darmstadt, den \today
\newpage
\phantomsection
\addcontentsline{toc}{section}{Abstract}
\let\oldabstractname\abstractname
\renewcommand{\abstractname}{Abstract}
\begin{abstract}
 
\end{abstract}
\newpage
\phantomsection
\addcontentsline{toc}{section}{Zusammenfassung}
\let\abstractname\oldabstractname
\begin{abstract}
 
\end{abstract}
\newpage
\phantomsection
\addcontentsline{toc}{section}{Inhaltsverzeichnis}
\tableofcontents
% Letzte Seite mit römischer Seitenzahl speichern
\cleardoublepage
\setcounter{savepage}{\arabic{page}}
\newpage
\cleardoublepage
\pagenumbering{arabic}
% Struktur basiert auf http://kleinmann.eit.h-da.de/99_ThesisInfos/
\section{Einführung}
Das Zentrum für Robotik hat den Industrieroboter TX60L mit der speziellen Schnittstelle \ac{LLI} beschafft. Mit dieser Schnittstelle ist es für die h\_da möglich, eigene Steuerungen und Regelungen zu realisieren. Kinematische Parameter und Parameter der Dynamik wurden ermittelt bzw. identifiziert\cite{D.X.Nodem.2015}. Für das Gelenk 1 wurde eine Einzelgelenkregelung realisiert (s. auch Versuch Labor Regelung von Roboterarmen). Ein Greifer wurde beschafft, der in das Robotersystem integriert werden kann.
\subsection{Was ist die Aufgabe?}
Für Gelenk 2 ist eine dezentrale Lageregelung in Kaskadenform in Anlehnung an die schon vorhandene Regelung für Gelenk 1 zu entwerfen, zu simulieren und zu realisieren. Dabei ist auch eine Eigengewichtskompensation vorzusehen und das Regelungsverhalten bei Lastmassenabwurf zu testen.
\subsection{Abzusehende Teilaufgaben}
\begin{itemize}
	\item Einarbeitung in Modellbeschreibung
	\item Erstellung des Simulationsmodells und Tests bei geeigneten Bewegungsvorgaben
	\item Einarbeitung in die Steuerungsumgebung mit LLI-Schnittstelle
	\item Montage des Greifers und Integration der Greiferansteuerung in die Steuerungsumgebung
	\item Implementierung der Steuerungsalgorithmen
	\item Experimentdurchführung und Vergleich mit der Simulation
	\item Dokumentation
\end{itemize}
\subsection{Was ist das Ziel dieser Arbeit?}

\subsection{Was war der Status Quo, bevor mit dieser Arbeit begonnen wurde?}

\subsubsection{Firma X hat etwas in der Vergangenheit auf diesem oder jenem Wege gemacht…}

\section{Problemlösung}

\subsection{Was sind die Optionen / Welche prinzipiellen Lösungen sind möglich?}

\subsection{Stand der Technik / Wie wird es woanders gemacht?}

\subsection{Was ist der in dieser Arbeit gewählte Ansatz?}

\subsection{Weshalb wurde dieser Ansatz gewählt / Was unterscheidet ihn von anderen / Was ist neu / Was ist bekannt?}

\subsection{Detaillierte Beschreibung der Problemlösung}

\section{Implementierung und Test}

\subsection{Wie wurde es implementiert}

\subsection{Wie wurde getestet}

\subsection{Warum wurde so getestet?}

\section{Validierung}

\subsection{Was sind die Ergebnisse}

\subsection{Vor- und Nachteile des entwickelten Systems}

\section{Zusammenfassung / Fazit / zukünftige Arbeit}

\subsection{Was wurde getan}

\subsection{Was muss noch getan werden / Wie kann das System verbessert werden?}

\newpage
\cleardoublepage
\pagenumbering{Roman}
% Seitenzahl wiederherstellen
\setcounter{page}{\thesavepage}
\clearpage
\phantomsection
\addcontentsline{toc}{section}{Literatur}
\printbibliography
\newpage
\renewcommand{\indexname}{Stichwortverzeichnis}
\printindex
\addcontentsline{toc}{section}{Stichwortverzeichnis}
\newpage
\section*{Abkürzungsverzeichnis}
\addcontentsline{toc}{section}{Abkürzungsverzeichnis}
\begin{acronym}[Bash] % längste Abkürzung
	\acro{LLI}{Low Level Interface}
\end{acronym}
\newpage
\phantomsection
\addcontentsline{toc}{section}{Abbildungsverzeichnis}
\listoffigures
\newpage
\phantomsection
\addcontentsline{toc}{section}{Tabellenverzeichnis}
\listoftables
\newpage
\section*{Anhang}
\addcontentsline{toc}{section}{Anhang}
\subsection{Protokoll}
\begin{itemize}
	\item 7.4.2016
		\begin{itemize}
			\item Projekt zugeteilt
			\item erste Unterlagen erhalten
			\item Vorschlag von Kommilitone, der zuvor bei Reis arbeitete: Quaderförmige Greifbacken mit gummierten Oberflächen
		\end{itemize}
	\item 8.4.2016
		\begin{itemize}
			\item Überlegungen zum Greifer: evtl. konkave Greifbacken
			\item Frage: Welche Form hat die Lastmasse?
			\item Wiedereinlesen in Kaskadenregelung
			\item Regler vereinfacht einen Teil der Strecke		
			\item Informieren über LLI
			\item Low Level Interface
			\item LLI gibt Zugriff auf Servos (Position control, Drehmoment control), IO und Vorwärtskinematik und Inverse Kinematik Solver
			\item Programmierung in C / C++
			\item Code für Motion Control \& Bahnplanung muss selber geschrieben werden
			\item Hardwarenah
			\item gut für Forschungszwecke (Universitäten und Forschungsinstitute)
			\item mit Dokumentation begonnen
		\end{itemize}
\end{itemize}
\end{document}